\section{秦：商鞅入秦与一场高代价的重组}

\eventref{WQ-0356-SHANGYANG}{商鞅变法}常被讲成“秦强大，所以商鞅正确”。这样会把过程倒过来。商鞅进入的秦，拥有关中腹地和向东扩展的潜力，却也需要把旧贵族关系、军功回报、土地利用和基层控制重新编排，才能和魏、楚、齐等强国竞争。

传统材料称商鞅为卫人，曾在魏活动，后得到秦孝公任用。这个跨国轨迹本身很重要：战国的游士并非只传播口才，也带着制度经验、军事知识和对别国强弱的观察。

\begin{event}{公元前350年}{商鞅第二次变法与秦都咸阳}{可靠纪年}{秦}\eventanchor{WQ-0350-QINREFORM}{战国·秦迁咸阳}
\term{这一年发生了什么}：传统叙事把迁都咸阳、行政区划与进一步制度调整联系在商鞅改革阶段。
\term{后来改变了什么}：改革从奖惩口号走向更稳定的行政空间。国家不是只在战场上变强，而是在户籍、县邑、道路、市场和法律的日常运转中变强。
\end{event}

\begin{sourcenote}[商鞅该怎样读]
\sourcebook{史记·商君列传}提供最有故事性的叙事；《商君书》的篇章、成书和作者归属不能简单等同于商鞅本人。军功爵、什伍连坐、农战和郡县化是理解秦制的重要线索，但“开阡陌”等具体制度的含义和实施方式仍有争议。睡虎地、里耶秦简展示的是约百年后已经成熟运行的法律与县级行政，能证明秦制的长期结果，不能把每一条成熟制度都直接归作商鞅个人的即时作品。
\end{sourcenote}
